\documentclass{beamer}

\usepackage[utf8]{inputenc}
\usepackage{graphicx}
\usepackage{listings}

\usepackage{tikz}
\usepackage{xcolor}
\usepackage{tikz-3dplot}
\usetikzlibrary{3d}
\usetikzlibrary{decorations.pathreplacing,calligraphy}
\newif\ifshowcellnumber
\showcellnumbertrue

\lstset{basicstyle=\ttfamily,breaklines=true}
\usepackage{fontspec}
\makeatletter
\def\verbatim@nolig@list{}
\makeatother
\setmonofont{JetBrains Mono}[Contextuals=Alternate]

\usepackage{amsmath}
\usepackage{namedtensor}
\usepackage{amsfonts}
\usepackage{adjustbox}
\usepackage{unicode-math}
\usepackage[pdf]{graphviz}

\title{AD in PyTorch}
\titlegraphic{\includegraphics[width=0.3\textwidth]{pytorch}}
\author{Breandan Considine}

\begin{document}
    \frame{\titlepage}


    \begin{frame}
        \frametitle{What is automatic differentiation?}
        To understand AD, you just need to remember two simple rules:
        \begin{align*}
            D(f + g) &= D(f) + D(g) \\
            D(f \cdot g) &= D(f) \cdot g + f \cdot D(g)
        \end{align*}
        We can think of AD as a \textit{linear map} between function spaces.

        \begin{align*}
            D(f + g) &= D(f) + D(g) \\
            D(f \cdot g) &= D(f) \cdot g + f \cdot D(g)
        \end{align*}
    \end{frame}


    \begin{frame}
        \frametitle{Higher order and higher rank AD}

        The \textit{gradient}, $\nabla: (\mathbb{R}^m\rightarrow\mathbb{R})\rightarrow\mathbb{R}^m$ maps a function $Q$ to:

        \begin{equation*}
            \nabla Q(q_1, \dots, q_m) = \left[ \frac{\partial Q}{\partial q_1}, \dots, \frac{\partial Q}{\partial q_m}\right]
        \end{equation*}

        The \textit{Jacobian}, $\mathcal{J}: (\mathbb{R}^m\rightarrow\mathbb{R}^n)\rightarrow\mathbb{R}^{n \times m}$ is a matrix of partials:

        \begin{equation*}
            \mathcal{J}\circ\mathbf{f} =
            \begin{bmatrix}
                \dfrac{\partial \mathbf{f}}{\partial x_1} & \cdots & \dfrac{\partial \mathbf{f}}{\partial x_m}
            \end{bmatrix} =
            \begin{bmatrix}
                \dfrac{\partial f_1}{\partial x_1} & \cdots & \dfrac{\partial f_1}{\partial x_m}\\
                \vdots & \ddots & \vdots\\
                \dfrac{\partial f_n}{\partial x_1} & \cdots & \dfrac{\partial f_n}{\partial x_m}
            \end{bmatrix} =
            \begin{bmatrix}
                \nabla f_1 \\
                \vdots \\
                \nabla f_m
            \end{bmatrix}
        \end{equation*}
    \end{frame}

    \begin{frame}[fragile]
        \frametitle{Gradients in PyTorch}

        Suppose we have a scalar-valued vector function, $f: \mathbb{R}^{n}\times\mathbb{R}^{n}\rightarrow\mathbb{R}$:

        \noindent\rule{\textwidth}{0.5pt}

        \begin{lstlisting}[language=Python]
x = torch.randn(2, requires_grad=True)
t = torch.tensor([2., 3.], requires_grad=True)
y = torch.randn(2)
f = sum((x*t - y)**2)**0.5

torch.autograd.grad(f, inputs=[x, t])
        \end{lstlisting}

        \noindent\rule{\textwidth}{0.5pt}
        \begin{lstlisting}

(tensor([-0.60,  2.85]), tensor([0.23, 1.47]))
        \end{lstlisting}

    \end{frame}

    \begin{frame}[fragile]
        \frametitle{Jacobians in PyTorch}
        \begin{lstlisting}[language=Python]
def jacobian(fun, x) -> torch.Tensor:
    x = x.detach().requires_grad_()
    y = fun(x)
    def vjp(v) -> torch.Tensor:
        return torch.autograd.grad(y, x, v)[0]

    vs = torch.eye(y.numel())\
              .view(y.numel(), *y.shape)
    result = vmap(vjp)(vs)
    return result.detach()

f = lambda x: x ** 3
jacobian(f, torch.ones(3))
        \end{lstlisting}
    \end{frame}

    \begin{frame}
        \frametitle{Higher order and higher rank AD}

        Suppose we have a function $P(X) = p_k \circ p_{k-1} \circ \cdots \circ p_1 \circ X$. The derivative of a linear composition can be expressed as a product:

        \begin{equation*}
            \label{eq:sfun_chain_rule}
            \frac{dP}{dp_1} = \frac{dp_k}{dp_{k-1}}\frac{dp_{k-1}}{dp_{k-2}}\dots\frac{dp_2}{dp_1}= {\displaystyle \prod_{i=1}^{k-1} \frac{dp_{i+1}}{dp_{i}}}
        \end{equation*}

        This also holds in higher dimensions, for example $\mathbf{P}_k: \mathbb{R}^m\rightarrow\mathbb{R}^n$:

        \begin{align*}
            \label{eq:vfun_chain_rule}
            \mathcal{J} \mathbf{P_k} = \displaystyle \prod_{i=1}^{k} \mathcal{J}p_i &= \underbrace{\bigg(\Big((\mathcal{J}p_k \mathcal{J}p_{k-1}) \dots \mathcal{J}p_2\Big) \mathcal{J}p_1\bigg)}_{\textit{Reverse mode, VJP, Pullback}} \\
            &= \underbrace{\bigg(\mathcal{J}p_k \Big(\mathcal{J}p_{k-1} \dots (\mathcal{J}p_2 \mathcal{J}p_1)\Big)\bigg)}_{\textit{Forward mode, JVP, Pushforward}}
        \end{align*}
    \end{frame}

    \begin{frame}
        \frametitle{Higher order and higher rank AD}

        The \textit{Hessian} $\mathbf{H}:(\mathbb{R}^m\rightarrow\mathbb{R})\rightarrow\mathbb{R}^{m\times m}$ maps scalar functions to $\partial^2$:

        \begin{equation*}
            \mathbf{H}(Q) = \begin{bmatrix}{\dfrac {\partial ^{2}Q}{\partial x_{1}^{2}}}
                                &{\dfrac {\partial ^{2}Q}{\partial x_{1}\,\partial x_{2}}}&\cdots &{\dfrac {\partial ^{2}Q}{\partial x_{1}\,\partial x_{m}}}\\[2.2ex]{\dfrac {\partial ^{2}Q}{\partial x_{2}\,\partial x_{1}}}&{\dfrac {\partial ^{2}Q}{\partial x_{2}^{2}}}&\cdots &{\dfrac {\partial ^{2}Q}{\partial x_{2}\,\partial x_{m}}}\\[2.2ex]\vdots &\vdots &\ddots &\vdots \\[2.2ex]{\dfrac {\partial ^{2}Q}{\partial x_{m}\,\partial x_{1}}}&{\dfrac {\partial ^{2}Q}{\partial x_{m}\,\partial x_{2}}}&\cdots &{\dfrac {\partial ^{2}Q}{\partial x_{m}^{2}}}
            \end{bmatrix}
        \end{equation*}

        The Hessian and Jacobian are related by $\mathbf{H}(Q)^\intercal = \mathcal{J} \circ \nabla Q$.
    \end{frame}

    \begin{frame}[fragile]
        \frametitle{Hessians in PyTorch}
        \begin{lstlisting}[language=Python]
def hessian(fun, x) -> torch.Tensor:
    def grad0(x: torch.Tensor):
        y = fun(x)
        assert y.dim() == 0
        return torch.autograd.grad(y, x, create_graph=True)[0]
    return jacobian(grad0, x)

g = lambda x: (x ** 3).sum()
hessian(g, torch.ones(3))
        \end{lstlisting}
    \end{frame}

    \begin{frame}[fragile]
        \frametitle{Static vs. Dynamic ADs: Representations}
        \begin{table}[H]
            \centering
            \begin{tabular}{lll}
                \textbf{Program} & & \textbf{Computation Graph} \\\\
%            \begin{tabular}[c]{@{}l@{}} $\hat y = θx + b$\\ $l = ||\hat y - y||_2$\end{tabular} &
                \begin{lstlisting}[basicstyle=\ttfamily\footnotesize, language=Python]
sum = 0
l = [0, 0, 0, 0]
for i in range(0, 4):
  l[i] += t[i] * x[i]
for i in range(0, 4):
  l[i] -= y[i] - b
for i in range(0, 4):
  l[i] *= l[i]
for i in range(0, 4):
  sum += l[i]
l = sqrt(sum)
                \end{lstlisting}                            & &
                \begin{adjustbox}{minipage={.4\textwidth}, height=0.22\textwidth, margin*=0cm 0cm 0cm 0.1cm}
                    \digraph[scale=0.1]{prograph}{
                        node[ fontname="Helvetica" fontsize=20 shape=Mrecord ];
                        edge[ fontname="Helvetica" fontsize=18 ];

                        graph ["concentrate"="true","rankdir"="LR","bgcolor"="transparent","margin"="0.0","compound"="true","nslimit"="20"]
                        "eeba8" ["label"="+"]
                        "a8416" ["label"="+"]
                        "4500f" ["label"="pow"]
                        "a67f9" ["label"="*"]
                        "0.5" ["label"="0.5"]
                        "f14a3" ["label"="*"]
                        "9c49d" ["label"="*"]
                        "59c48" ["label"="+"]
                        "980bd" ["label"="+"]
                        "8f532" ["label"="+"]
                        "1a609" ["label"="+"]
                        "e58c4" ["label"="+"]
                        "23f5b" ["label"="+"]
                        "d829b" ["label"="+"]
                        "y₀" ["label"="y₀"]
                        "2783d" ["label"="+"]
                        "8bd47" ["label"="+"]
                        "y₂" ["label"="y₂"]
                        "517e6" ["label"="+"]
                        "8caa0" ["label"="+"]
                        "y₁" ["label"="y₁"]
                        "7da0d" ["label"="+"]
                        "b12cb" ["label"="+"]
                        "b" ["label"="b"]
                        "f8941" ["label"="+"]
                        "3eecd" ["label"="+"]
                        "2e83a" ["label"="+"]
                        "b59fd" ["label"="+"]
                        "dae83" ["label"="+"]
                        "b11ba" ["label"="*"]
                        "3bb89" ["label"="*"]
                        "b2454" ["label"="*"]
                        "7bed4" ["label"="*"]
                        "39644" ["label"="*"]
                        "12c32" ["label"="*"]
                        "d58d1" ["label"="*"]
                        "6c64d" ["label"="*"]
                        "fb0f0" ["label"="*"]
                        "6c2be" ["label"="*"]
                        "57fd4" ["label"="*"]
                        "a9bc3" ["label"="*"]
                        "x₀" ["label"="x₀"]
                        "θ₀" ["label"="θ₀"]
                        "x₂" ["label"="x₂"]
                        "θ₁" ["label"="θ₁"]
                        "x₂" ["label"="x₂"]
                        "x₄" ["label"="x₄"]
                        "x₁" ["label"="x₁"]
                        "x₃" ["label"="x₃"]
                        "eeba8" -> "a8416"
                        "a8416" -> "4500f"
                        "a67f9" -> "a8416"
                        "0.5" -> "4500f"
                        "f14a3" -> "eeba8"
                        "9c49d" -> "eeba8"
                        "59c48" -> "a67f9"
                        "980bd" -> "a67f9"
                        "8f532" -> "f14a3"
                        "1a609" -> "f14a3"
                        "e58c4" -> "9c49d"
                        "23f5b" -> "9c49d"
                        "d829b" -> "59c48"
                        "y₀" -> "59c48"
                        "y₀" -> "980bd"
                        "2783d" -> "980bd"
                        "8bd47" -> "8f532"
                        "y₂" -> "8f532"
                        "y₂" -> "1a609"
                        "517e6" -> "1a609"
                        "8caa0" -> "e58c4"
                        "y₁" -> "e58c4"
                        "y₁" -> "23f5b"
                        "7da0d" -> "23f5b"
                        "b12cb" -> "d829b"
                        "b" -> "d829b"
                        "b" -> "2783d"
                        "b" -> "8bd47"
                        "b" -> "517e6"
                        "b" -> "8caa0"
                        "b" -> "7da0d"
                        "f8941" -> "2783d"
                        "3eecd" -> "8bd47"
                        "2e83a" -> "517e6"
                        "b59fd" -> "8caa0"
                        "dae83" -> "7da0d"
                        "b11ba" -> "b12cb"
                        "3bb89" -> "b12cb"
                        "b2454" -> "f8941"
                        "7bed4" -> "f8941"
                        "39644" -> "3eecd"
                        "12c32" -> "3eecd"
                        "d58d1" -> "2e83a"
                        "6c64d" -> "2e83a"
                        "fb0f0" -> "b59fd"
                        "6c2be" -> "b59fd"
                        "57fd4" -> "dae83"
                        "a9bc3" -> "dae83"
                        "x₀" -> "b11ba"
                        "x₀" -> "b2454"
                        "θ₀" -> "b11ba"
                        "θ₀" -> "b2454"
                        "θ₀" -> "39644"
                        "θ₀" -> "d58d1"
                        "θ₀" -> "fb0f0"
                        "θ₀" -> "57fd4"
                        "x₂" -> "3bb89"
                        "x₂" -> "7bed4"
                        "θ₁" -> "3bb89"
                        "θ₁" -> "7bed4"
                        "θ₁" -> "12c32"
                        "θ₁" -> "6c64d"
                        "θ₁" -> "6c2be"
                        "θ₁" -> "a9bc3"
                        "x₂" -> "39644"
                        "x₂" -> "d58d1"
                        "x₄" -> "12c32"
                        "x₄" -> "6c64d"
                        "x₁" -> "fb0f0"
                        "x₁" -> "57fd4"
                        "x₃" -> "6c2be"
                        "x₃" -> "a9bc3"
                    }
                \end{adjustbox}
            \end{tabular}
        \end{table}
    \end{frame}

%    \begin{frame}
%      \frametitle{AD in PyTorch}
%      \begin{enumerate}
%        \item Introduction to AD
%        \item Implicit vs.
%      \end{enumerate}
%    \end{frame}
%
%    \begin{frame}
%      \frametitle{Stateful vs. stateless computation}
%      \begin{enumerate}
%        \item Introduction to AD
%        \item Implicit vs.
%      \end{enumerate}
%    \end{frame}

    \begin{frame}
        \frametitle{What is a tensor?}

        \definecolor{R}{RGB}{202,65,55}
        \definecolor{G}{RGB}{151,216,56}
        \definecolor{B}{RGB}{0,0,0}
        \definecolor{W}{RGB}{255,255,255}
        \definecolor{X}{RGB}{65,65,65}

        \newcommand{\TikZRubikFaceLeft}[9]{\def\myarrayL{#1,#2,#3,#4,#5,#6,#7,#8,#9}}
        \newcommand{\TikZRubikFaceRight}[9]{\def\myarrayR{#1,#2,#3,#4,#5,#6,#7,#8,#9}}
        \newcommand{\TikZRubikFaceTop}[9]{\def\myarrayT{#1,#2,#3,#4,#5,#6,#7,#8,#9}}
        \newcommand{\BuildArray}{\foreach \X [count=\Y] in \myarrayL%
        {\ifnum\Y=1%
        \xdef\myarray{"\X"}%
        \else%
        \xdef\myarray{\myarray,"\X"}%
        \fi}%
        \foreach \X in \myarrayR%
        {\xdef\myarray{\myarray,"\X"}}%
        \foreach \X in \myarrayT%
        {\xdef\myarray{\myarray,"\X"}}%
        \xdef\myarray{{\myarray}}%
        }
        \TikZRubikFaceLeft
        {X}{W}{W}
        {W}{X}{X}
        {X}{W}{W}
        \TikZRubikFaceRight
        {W}{X}{W}
        {X}{W}{X}
        {W}{X}{W}
        \TikZRubikFaceTop
        {X}{W}{X}
        {W}{W}{X}
        {W}{X}{W}
        \BuildArray
        \pgfmathsetmacro\radius{0.1}
        \tdplotsetmaincoords{55}{135}

        \showcellnumberfalse

        \bgroup
        \def\arraystretch{1.2}
        \begin{table}[H]
            \centering
            \begin{tabular}{cc}
                \textbf{Rank-2} & \textbf{Rank-3} \\\\
                \begin{adjustbox}{minipage={.49\textwidth}, margin*=-0.5cm 0cm -0.5cm 0cm}
                    \begin{equation*}
                        \begin{bmatrix}{1}
                            &{0}&\cdots &{1}\\{0}&{1}&\cdots &{0}\\\vdots &\vdots &\ddots &\vdots \\{1}&{0}&\cdots &{1}
                        \end{bmatrix}
                    \end{equation*}
                \end{adjustbox} &
                \begin{adjustbox}{minipage={.48\textwidth}}
                    \begin{tikzpicture}
                        \clip (-3,-2.5) rectangle (3,2.5);
                        \begin{scope}[tdplot_main_coords]
                            \filldraw [canvas is yz plane at x=1.5] (-1.5,-1.5) rectangle (1.5,1.5);
                            \filldraw [canvas is xz plane at y=1.5] (-1.5,-1.5) rectangle (1.5,1.5);
                            \filldraw [canvas is yx plane at z=1.5] (-1.5,-1.5) rectangle (1.5,1.5);
                            \foreach \X [count=\XX starting from 0] in {-1.5,-0.5,0.5}{
                                \foreach \Y [count=\YY starting from 0] in {-1.5,-0.5,0.5}{
                                    \pgfmathtruncatemacro{\Z}{\XX+3*(2-\YY)}
                                    \pgfmathsetmacro{\mycolor}{\myarray[\Z]}
                                    \draw [thick,canvas is yz plane at x=1.5,shift={(\X,\Y)},fill=\mycolor] (0.5,0) -- ({1-\radius},0) arc (-90:0:\radius) -- (1,{1-\radius}) arc (0:90:\radius) -- (\radius,1) arc (90:180:\radius) -- (0,\radius) arc (180:270:\radius) -- cycle;
                                    \ifshowcellnumber
                                    \node[canvas is yz plane at x=1.5,shift={(\X+0.5,\Y+0.5)}] {\Z};
                                    \fi
                                    \pgfmathtruncatemacro{\Z}{2-\XX+3*(2-\YY)+9}
                                    \pgfmathsetmacro{\mycolor}{\myarray[\Z]}
                                    \draw [thick,canvas is xz plane at y=1.5,shift={(\X,\Y)},fill=\mycolor] (0.5,0) -- ({1-\radius},0) arc (-90:0:\radius) -- (1,{1-\radius}) arc (0:90:\radius) -- (\radius,1) arc (90:180:\radius) -- (0,\radius) arc (180:270:\radius) -- cycle;
                                    \ifshowcellnumber
                                    \node[canvas is xz plane at y=1.5,shift={(\X+0.5,\Y+0.5)},xscale=-1] {\Z};
                                    \fi
                                    \pgfmathtruncatemacro{\Z}{2-\YY+3*\XX+18}
                                    \pgfmathsetmacro{\mycolor}{\myarray[\Z]}
                                    \draw [thick,canvas is yx plane at z=1.5,shift={(\X,\Y)},fill=\mycolor] (0.5,0) -- ({1-\radius},0) arc (-90:0:\radius) -- (1,{1-\radius}) arc (0:90:\radius) -- (\radius,1) arc (90:180:\radius) -- (0,\radius) arc (180:270:\radius) -- cycle;
                                    \ifshowcellnumber
                                    \node[canvas is yx plane at z=1.5,shift={(\X+0.5,\Y+0.5)},xscale=-1,rotate=-90] {\Z};
                                    \fi
                                }
                            }

                            \draw [decorate,decoration={calligraphic brace,amplitude=10pt,mirror},yshift=0pt, line width=1.25pt]
                            (3,0) -- (3,3) node [black,midway,xshift=-8pt, yshift=-14pt] {\footnotesize $\mathcal{E}$};
                            \draw [decorate,decoration={calligraphic brace,amplitude=10pt},yshift=0pt, line width=1.25pt]
                            (3,0) -- (0,-3) node [black,midway,xshift=-16pt, yshift=0pt] {\footnotesize $\mathcal{F}$};
                            \draw [decorate,decoration={calligraphic brace,amplitude=10pt},yshift=0pt, line width=1.25pt]
                            (0,-3) -- (-3,-3) node [black,midway,xshift=-8pt, yshift=14pt] {\footnotesize $\mathcal{G}$};
                        \end{scope}
                    \end{tikzpicture}
                \end{adjustbox}
            \end{tabular}
        \end{table}
        \egroup
    \end{frame}

    \begin{frame}
        \frametitle{Tensor contraction}
        \begin{table}[H]
            \begin{tabular}{lll}
                \textbf{Notation} & & \textbf{Code}                 \\\\
                $A \ndot{foo} B$ & & \texttt{A.mv(B)} \\

                $\nsum{foo} C$ & & \texttt{namedtensor.sum('C')}
            \end{tabular}
        \end{table}
    \end{frame}

    \begin{frame}[fragile]
        \frametitle{Checking matrix multiplication}

        Suppose we have two tensors, $A: \mathbb{R}^{\textcolor{blue}{x}\times \textcolor{orange}{y} \times \cdots}$ and $B: \mathbb{R}^{\textcolor{orange}{y}\times \textcolor{purple}{z} \times \cdots}$

        Then \texttt{C = A @ B} has type $C: \mathbb{R}^{\textcolor{blue}{x}\times \textcolor{purple}{z} \times \cdots}$.

        \noindent\rule{\textwidth}{0.5pt}

        \begin{lstlisting}[escapeinside={<@}{@>}]

state = torch.ones(<@\textcolor{blue}{9}@>, <@\textcolor{orange}{5}@>, names=('<@\textcolor{blue}{batch}@>', '<@\textcolor{orange}{D}@>'))
trans = torch.randn(<@\textcolor{orange}{5}@>, <@\textcolor{purple}{5}@>, names=('<@\textcolor{orange}{in}@>', '<@\textcolor{purple}{out}@>'))
next_state = state @ trans
print(next_state.names)
        \end{lstlisting}
        \noindent\rule{\textwidth}{0.5pt}
        \begin{lstlisting}[escapeinside={<@}{@>}]

('<@\textcolor{blue}{batch}@>', '<@\textcolor{purple}{out}@>')
        \end{lstlisting}
    \end{frame}

    \begin{frame}[fragile]
        \frametitle{Runtime type checking: name mismatch}
        \begin{lstlisting}[language=Python]
x = torch.ones(3, names=('X',))
y = torch.ones(3)
z = torch.ones(3, names=('Z',))
xpz = x + z
        \end{lstlisting}
        \noindent\rule{\textwidth}{0.5pt}
        \begin{lstlisting}[escapeinside={<@}{@>}]
<@\textcolor{red}{----RuntimeError}@>
Traceback (most recent call last)[...]
      <@\textcolor{green}{2}@> y = torch.ones(3)
      <@\textcolor{green}{3}@> z = torch.ones(3, names=('Z',))
<@\textcolor{green}{----> 4}@> xpz = x + z
<@\textcolor{red}{RuntimeError}@>: Error attempting to broadcast
dims ['X'] and dims ['Z']: dim 'X' and dim 'Z'
are at the same position from the right
        \end{lstlisting}
    \end{frame}


\end{document}